% GENERATED FILE. Edit canonical lessons, then run npm run generate:books.
% canonical-source-hash: fnv1a64:ca28bbe1d54f8b25
% canonical-lessons: SA-C543-sveccha, SA-C543-vivasah, SA-C543-hathah, SA-C543-agrahah, SA-C543-svayam

\chapter{Free to Choose}
\label{ch:sa-free-to-choose}
\hlchaptermodality{543}

\begin{chapteropening}
\textbf{By the end of this chapter:} \emph{I can say whether someone acts freely, stubbornly or helplessly.}

Complete the last lesson of chapter 543: I can say whether someone acts freely, stubbornly or helplessly.
\end{chapteropening}

\section[\texorpdfstring{\sk{स्वेच्छा} (svecchā)}{स्वेच्छा (svecchā)}]{\sk{स्वेच्छा} (svecchā) — free will}

\label{lesson:SA-C543-sveccha}



\begin{quote}
Before the new one: say the Sanskrit for pleasing, then the Sanskrit for neutral.
\end{quote}

\subsection*{You'll want to know: \sk{स्वेच्छा}}
\textbf{\sk{स्वेच्छा}} — \emph{svecchā} — \textquotedblleft{}free will\textquotedblright{}.

\begin{quote}
\sk{यदि} \sk{सः} \sk{स्वेच्छया} \sk{आगच्छति}, \sk{तर्हि} \sk{वयं} \sk{प्रसन्नाः} \sk{स्मः।}
\end{quote}

\begin{quote}
\emph{yadi saḥ svecchayā āgacchati, tarhi vayaṃ prasannāḥ smaḥ.}
\end{quote}

\begin{quote}
If he comes of his own free will, we are glad.
\end{quote}

A feminine noun.

\begin{cousinweb}
\textbf{\sk{स्वेच्छा}} (\emph{svecchā}) is built from \emph{sva}, \textquotedblleft{}own\textquotedblright{}, and \emph{icchā}, \textquotedblleft{}a wish\textquotedblright{}, from the root √iṣ, \textquotedblleft{}to wish\textquotedblright{}, which also gives \emph{icchati}; MW gives \emph{svecchā} as one's own wish, free will (MW).
\end{cousinweb}

\begin{grammarlens}[title={What you've built}]
One more word for saying what happens if.
\end{grammarlens}

\subsection*{Your turn}
\begin{itemize}
\raggedright
  \item \emph{Say it:} \emph{svecchā}
  \item \emph{Say it:} \emph{svecchā}, once more
  \item \emph{Say it:} \emph{yadi saḥ svecchayā āgacchati, tarhi vayaṃ prasannāḥ smaḥ.} — If he comes of his own free will, we are glad.
  \item \emph{Recall:} say the Sanskrit for pleasing, then the Sanskrit for neutral, then say \emph{svecchā} again
\end{itemize}

\begin{tcolorbox}[breakable,colback=teal!4,colframe=teal!35!black,title={Before you move on}]
What does \sk{स्वेच्छा} mean? (\textquotedblleft{}Free will\textquotedblright{}.) Say it once more.
\end{tcolorbox}
\section[\texorpdfstring{\sk{विवशः} (vivaśaḥ)}{विवशः (vivaśaḥ)}]{\sk{विवशः} (vivaśaḥ) — helpless}

\label{lesson:SA-C543-vivasah}



\begin{quote}
Before the new one: say the Sanskrit for neutral, then the Sanskrit for free will.
\end{quote}

\subsection*{You'll want to know: \sk{विवशः}}
\textbf{\sk{विवशः}} — \emph{vivaśaḥ} — \textquotedblleft{}helpless\textquotedblright{}.

\begin{quote}
\sk{यदि} \sk{सः} \sk{विवशः} \sk{अस्ति}, \sk{तर्हि} \sk{तस्य} \sk{साहाय्यं} \sk{करोतु।}
\end{quote}

\begin{quote}
\emph{yadi saḥ vivaśaḥ asti, tarhi tasya sāhāyyaṃ karotu.}
\end{quote}

\begin{quote}
If he is helpless, please help him.
\end{quote}

\begin{cousinweb}
\textbf{\sk{विवशः}} (\emph{vivaśaḥ}) is built on \emph{vi}, \textquotedblleft{}without\textquotedblright{}, and \emph{vaśa}, \textquotedblleft{}will, control\textquotedblright{}; MW gives \emph{vivaśa} as deprived of will, powerless, helpless (MW).
\end{cousinweb}

\begin{grammarlens}[title={What you've built}]
One more word for saying what happens if.
\end{grammarlens}

\subsection*{Your turn}
\begin{itemize}
\raggedright
  \item \emph{Say it:} \emph{vivaśaḥ}
  \item \emph{Say it:} \emph{vivaśaḥ}, once more
  \item \emph{Say it:} \emph{yadi saḥ vivaśaḥ asti, tarhi tasya sāhāyyaṃ karotu.} — If he is helpless, please help him.
  \item \emph{Recall:} say the Sanskrit for neutral, then the Sanskrit for free will, then say \emph{vivaśaḥ} again
\end{itemize}

\begin{tcolorbox}[breakable,colback=teal!4,colframe=teal!35!black,title={Before you move on}]
What does \sk{विवशः} mean? (\textquotedblleft{}Helpless\textquotedblright{}.) Say it once more.
\end{tcolorbox}
\section[\texorpdfstring{\sk{हठः} (haṭhaḥ)}{हठः (haṭhaḥ)}]{\sk{हठः} (haṭhaḥ) — obstinacy}

\label{lesson:SA-C543-hathah}



\begin{quote}
Before the new one: say the Sanskrit for free will, then the Sanskrit for helpless.
\end{quote}

\subsection*{You'll want to know: \sk{हठः}}
\textbf{\sk{हठः}} — \emph{haṭhaḥ} — \textquotedblleft{}obstinacy\textquotedblright{}.

\begin{quote}
\sk{यदि} \sk{बालः} \sk{हठं} \sk{करोति}, \sk{तर्हि} \sk{माता} \sk{कुपिता} \sk{भवति।}
\end{quote}

\begin{quote}
\emph{yadi bālaḥ haṭhaṃ karoti, tarhi mātā kupitā bhavati.}
\end{quote}

\begin{quote}
If the child is stubborn, my mother gets angry.
\end{quote}

A masculine noun.

\begin{cousinweb}
\textbf{\sk{हठः}} (\emph{haṭhaḥ}) is the word MW gives for violence, force, and obstinacy, pertinacity; MW gives no derivation, so its origin is uncertain (MW).
\end{cousinweb}

\begin{grammarlens}[title={What you've built}]
One more word for saying what happens if.
\end{grammarlens}

\subsection*{Your turn}
\begin{itemize}
\raggedright
  \item \emph{Say it:} \emph{haṭhaḥ}
  \item \emph{Say it:} \emph{haṭhaḥ}, once more
  \item \emph{Say it:} \emph{yadi bālaḥ haṭhaṃ karoti, tarhi mātā kupitā bhavati.} — If the child is stubborn, my mother gets angry.
  \item \emph{Recall:} say the Sanskrit for free will, then the Sanskrit for helpless, then say \emph{haṭhaḥ} again
\end{itemize}

\begin{tcolorbox}[breakable,colback=teal!4,colframe=teal!35!black,title={Before you move on}]
What does \sk{हठः} mean? (\textquotedblleft{}Obstinacy\textquotedblright{}.) Say it once more.
\end{tcolorbox}
\section[\texorpdfstring{\sk{आग्रहः} (āgrahaḥ)}{आग्रहः (āgrahaḥ)}]{\sk{आग्रहः} (āgrahaḥ) — insistence}

\label{lesson:SA-C543-agrahah}



\begin{quote}
Before the new one: say the Sanskrit for helpless, then the Sanskrit for obstinacy.
\end{quote}

\subsection*{You'll want to know: \sk{आग्रहः}}
\textbf{\sk{आग्रहः}} — \emph{āgrahaḥ} — \textquotedblleft{}insistence\textquotedblright{}.

\begin{quote}
\sk{यदि} \sk{भवतः} \sk{आग्रहः} \sk{अस्ति}, \sk{तर्हि} \sk{अहम्} \sk{आगच्छामि।}
\end{quote}

\begin{quote}
\emph{yadi bhavataḥ āgrahaḥ asti, tarhi aham āgacchāmi.}
\end{quote}

\begin{quote}
If you insist, I will come.
\end{quote}

A masculine noun.

\begin{cousinweb}
\textbf{\sk{आग्रहः}} (\emph{āgrahaḥ}) is built on \emph{ā} and the root √grah, \textquotedblleft{}to seize\textquotedblright{}; MW gives \emph{āgraha} as insisting on, a strong inclination, obstinacy (MW).
\end{cousinweb}

\begin{grammarlens}[title={What you've built}]
One more word for saying what happens if.
\end{grammarlens}

\subsection*{Your turn}
\begin{itemize}
\raggedright
  \item \emph{Say it:} \emph{āgrahaḥ}
  \item \emph{Say it:} \emph{āgrahaḥ}, once more
  \item \emph{Say it:} \emph{yadi bhavataḥ āgrahaḥ asti, tarhi aham āgacchāmi.} — If you insist, I will come.
  \item \emph{Recall:} say the Sanskrit for helpless, then the Sanskrit for obstinacy, then say \emph{āgrahaḥ} again
\end{itemize}

\begin{tcolorbox}[breakable,colback=teal!4,colframe=teal!35!black,title={Before you move on}]
What does \sk{आग्रहः} mean? (\textquotedblleft{}Insistence\textquotedblright{}.) Say it once more.
\end{tcolorbox}
\section[\texorpdfstring{\sk{स्वयम्} (svayam)}{स्वयम् (svayam)}]{\sk{स्वयम्} (svayam) — by oneself}

\label{lesson:SA-C543-svayam}



\begin{quote}
Before the new one: say the Sanskrit for obstinacy, then the Sanskrit for insistence.
\end{quote}

\subsection*{You'll want to know: \sk{स्वयम्}}
\textbf{\sk{स्वयम्}} — \emph{svayam} — \textquotedblleft{}by oneself\textquotedblright{}.

\begin{quote}
\sk{यदि} \sk{त्वं} \sk{स्वयं} \sk{पठसि}, \sk{तर्हि} \sk{त्वं} \sk{शीघ्रं} \sk{शिक्षसे।}
\end{quote}

\begin{quote}
\emph{yadi tvaṃ svayaṃ paṭhasi, tarhi tvaṃ śīghraṃ śikṣase.}
\end{quote}

\begin{quote}
If you study by yourself, you learn fast.
\end{quote}

\begin{cousinweb}
\textbf{\sk{स्वयम्}} (\emph{svayam}) probably comes from an old nominative of \emph{sva}, \textquotedblleft{}own\textquotedblright{}, as MW takes it; MW gives \emph{svayam} as self, by oneself, of one's own accord (MW).
\end{cousinweb}

\begin{grammarlens}[title={What you've built}]
One more word for saying what happens if.
\end{grammarlens}

\subsection*{Your turn}
\begin{itemize}
\raggedright
  \item \emph{Say it:} \emph{svayam}
  \item \emph{Say it:} \emph{svayam}, once more
  \item \emph{Say it:} \emph{yadi tvaṃ svayaṃ paṭhasi, tarhi tvaṃ śīghraṃ śikṣase.} — If you study by yourself, you learn fast.
  \item \emph{Recall:} say the Sanskrit for obstinacy, then the Sanskrit for insistence, then say \emph{svayam} again
\end{itemize}

\begin{tcolorbox}[breakable,colback=teal!4,colframe=teal!35!black,title={Before you move on}]
What does \sk{स्वयम्} mean? (\textquotedblleft{}By oneself\textquotedblright{}.) Say it once more.
\end{tcolorbox}
